% Capítulo 6 (Resultados) — Slide 2 — Experimento 2: escopo
\begin{frame}{Experimento 2: expansão na Bahia}
  \begin{itemize}
    \item 145 séries (loja, produto), 3 filiais, 6 segmentos comerciais
    \item 12 políticas $\times$ 5 replicações independentes
    \item 8.700 pontos de avaliação no total
    \item 71\% das séries classificadas como \textit{Lumpy} no recorte avaliado
  \end{itemize}
  \vfill
  \small Cinco replicações por série habilitam testes estatísticos formais,
  o que o piloto da Paraíba (1 replicação) não permitia.
  \note{
    O Experimento 2 amplia a escala em relação ao piloto: 145 combinações
    loja-produto do estado da Bahia, com cinco replicações independentes por
    política, totalizando 8.700 pontos de avaliação. Essa escala é o que
    habilita os testes estatísticos formais -- Wilcoxon, Friedman, Cohen's
    d -- que serão apresentados mais adiante. A próxima imagem mostra como
    essas 145 séries se distribuem entre as filiais e segmentos comerciais
    da rede.
  }
\end{frame}
